\documentclass[11pt]{article}
\usepackage[margin=1in]{geometry}
\usepackage{amsmath,amssymb,amsthm}
\usepackage{xurl}
\usepackage{hyperref}
\sloppy
\let\oldus\_
\renewcommand{\_}{\oldus\allowbreak}

\newtheorem{theorem}{Theorem}
\newtheorem{lemma}[theorem]{Lemma}
\theoremstyle{remark}
\newtheorem{remark}[theorem]{Remark}
\newcommand{\Z}{\mathbb{Z}}
\newcommand{\GL}{\mathrm{GL}}

\title{Conjecture 00000005600:\\ same determinant, non-isomorphic echelon forms,\\
separated by a noncommutative unimodular pair}
\author{Jackmeson1 (AI-assisted)}
\date{2026-10-04}

\begin{document}
\maketitle

\section*{Statement}
\emph{English.} Definition: The determinant and the echelon form are two layers.
Conjecture: There exist two matrices with the same determinant but non-isomorphic
echelon forms, and the separation is realized by an explicit noncommutative pair of
unimodular transformations.
(The Chinese text says the same: the two matrices have equal determinants while their
echelon forms are not isomorphic, and the separation is realized by an explicit
noncommutative pair of unimodular transformations.)

\textbf{Verdict: proved.} The statement is existential; we give explicit integer
matrices and an explicit pair in $\GL_2(\Z)$, and the whole statement below is a single
Lean theorem, \texttt{Conjecture5600.conjecture\_5600}, checked against Mathlib.

\section{Definitions and reading}
All matrices are $2\times 2$ integer matrices. \emph{Unimodular transformations} are the
elements of $\GL_2(\Z)$ (integer matrices with integer inverse, equivalently
$\det=\pm1$); in Lean this is Mathlib's \texttt{GL (Fin 2) Z}. A unimodular row
transformation sends $A$ to $WA$, $W\in\GL_2(\Z)$; $A$ and $H$ are \emph{row equivalent}
if $H=WA$ for some such $W$.

\emph{Echelon form.} Over $\Z$ the echelon form of a nonsingular matrix is its row
Hermite normal form: $H=WA$ with $W\in\GL_2(\Z)$, $H$ upper triangular, positive
diagonal (pivot) entries, and every entry above a pivot in $[0,\text{pivot})$. This is the
Lean predicate \texttt{IsHermiteNormalForm} (stated for all $n\times n$).

\emph{Isomorphism type of an echelon form.} The row lattice $L(A)=\{wA : w\in\Z^2\}$
(Mathlib: \texttt{LinearMap.range A.vecMulLinear}, proved equal to the $\Z$-span of the
rows) is unchanged by row transformations, so all echelon forms of $A$ have the same
row lattice and the same cokernel $\operatorname{coker}A=\Z^2/L(A)$. We say two echelon
forms are \emph{non-isomorphic} in the strongest sense: their cokernels are not
isomorphic as abelian groups. This implies that they are different matrices, that the
matrices are not row equivalent, and that their Smith forms differ.

\emph{The separating pair.} We prove two readings, with the same explicit pair
\[
U=\begin{pmatrix}1&1\\0&1\end{pmatrix},\qquad
V=\begin{pmatrix}1&0\\1&1\end{pmatrix},\qquad
UV=\begin{pmatrix}2&1\\1&1\end{pmatrix}\neq
VU=\begin{pmatrix}1&1\\1&2\end{pmatrix},
\]
the two standard elementary generators of $\mathrm{SL}_2(\Z)$, with $\det U=\det V=1$.
\begin{itemize}
\item[(A)] \emph{The pair carries the two matrices to their separated echelon forms:}
$UM$ and $VN$ are the Hermite forms of $M$ and $N$, $\det M=\det N$, and the echelon
forms are non-isomorphic.
\item[(B)] \emph{The separation comes from the noncommutativity itself:} for one base
matrix $X$, the matrices $XUV$ and $XVU$ (the pair applied in the two orders) have equal
determinants, since $\det(UV)=\det(VU)$, but distinct echelon forms.
\end{itemize}

\section{Theorem and proof}
Let
\[
M=\begin{pmatrix}1&-4\\0&4\end{pmatrix},\quad
N=\begin{pmatrix}2&0\\-2&2\end{pmatrix},\quad
X=\begin{pmatrix}1&0\\0&2\end{pmatrix}.
\]

\begin{theorem}\label{thm:main}
\begin{enumerate}
\item $\det M=\det N=4$, $UM=\operatorname{diag}(1,4)$ and $VN=\operatorname{diag}(2,2)$
are in Hermite normal form, and $UM\neq VN$.
\item For all matrices $H_1$ row equivalent to $M$ and $H_2$ row equivalent to $N$
(in particular for all echelon forms of $M$ and $N$, under any convention),
$\operatorname{coker}H_1\not\cong\operatorname{coker}H_2$ and $H_1\neq H_2$.
\item There are no $P,Q\in\GL_2(\Z)$ with $PNQ=M$; so $M$ and $N$ are not even
equivalent under two-sided unimodular transformations, and their Smith forms differ.
\item $\det(XUV)=\det(XVU)=2$; $XUV$ has Hermite form $\operatorname{diag}(2,1)$ and
$XVU$ has Hermite form $\begin{pmatrix}1&1\\0&2\end{pmatrix}$; and every matrix row
equivalent to $XUV$ differs from every matrix row equivalent to $XVU$.
\end{enumerate}
\end{theorem}

\begin{proof}
(1) Direct multiplication: $\det M=4-0=4$, $\det N=4-0=4$,
$UM=\begin{pmatrix}1&0\\0&4\end{pmatrix}$, $VN=\begin{pmatrix}2&0\\0&2\end{pmatrix}$.
Both are upper triangular with positive diagonal and zero above-pivot entries, so they
are Hermite forms; their $(1,1)$ entries $1\neq 2$ differ.

(2) If $H=WA$ with $W\in\GL_2(\Z)$, then $wH=(wW)A$ and $wA=(wW^{-1})H$, so
$L(H)=L(A)$. Now $L(M)=\{(a,\,-4a+4b)\}=\Z\times 4\Z$ and
$L(N)=\{(2a-2b,\,2b)\}=2\Z\times 2\Z$. Every class $y$ in
$\Z^2/L(N)$ satisfies $y+y=0$, because $2v\in 2\Z\times 2\Z$ for every $v$. In
$\Z^2/L(M)$ the class $x$ of $(0,1)$ satisfies $x+x\neq 0$, because $(0,2)\notin\Z\times
4\Z$. An additive isomorphism $e$ would give $e(x+x)=e(x)+e(x)=0$, hence $x+x=0$, a
contradiction. (The cokernels are $\Z/4$ and $(\Z/2)^2$.) Equal matrices have equal
cokernels, so also $H_1\neq H_2$.

(3) All entries of $N$ are even, so all entries of $PNQ$ are even for all integer
$P,Q$; explicitly $(PNQ)_{11}=2\big((P_{11}-P_{12})Q_{11}+P_{12}Q_{21}\big)$, while
$M_{11}=1$.

(4) $XUV=\begin{pmatrix}2&1\\2&2\end{pmatrix}$ and
$XVU=\begin{pmatrix}1&1\\2&4\end{pmatrix}$, both of determinant $2$. With
$W_1=\begin{pmatrix}2&-1\\-1&1\end{pmatrix}$ and
$W_2=\begin{pmatrix}1&0\\-2&1\end{pmatrix}$ (determinant $1$) we get
$W_1XUV=\operatorname{diag}(2,1)$ and $W_2XVU=\begin{pmatrix}1&1\\0&2\end{pmatrix}$, both
Hermite forms. Every vector of $L(XUV)=\{(2a+2b,\,a+2b)\}$ has even first coordinate,
while $(1,1)$ is a row of $XVU$. By the invariance of $L$ under row transformations,
a matrix row equivalent to both would give $L(XUV)=L(XVU)$, which is false.
\end{proof}

\begin{remark}
In reading (B) both matrices are obtained from the same $X$ by unimodular column
operations, so they are two-sided equivalent and their cokernels are necessarily
isomorphic for every choice of $X$, $U$, $V$. In that reading, the layer that the
determinant fails to see is the row-equivalence class, which the echelon form
classifies; part (4) proves that these classes differ. Part (2) gives the cokernel
form of the separation for reading (A). The conjecture is existential, and the
theorem satisfies both readings with the same explicit noncommutative pair.
\end{remark}

\section{Lean formalization}
File \texttt{lean/Conjecture5600/Basic.lean} (Lean 4.33.1, Mathlib \texttt{v4.33.1}; the
root file \texttt{Conjecture5600.lean} only imports it). Main theorem:
\begin{verbatim}
theorem conjecture_5600 :
  exists (A B : Mat) (U V : GL (Fin 2) Int),
    (U : Mat).det = 1 /\ (V : Mat).det = 1 /\ U * V != V * U /\
    A.det = B.det /\
    IsHermiteNormalForm ((U : Mat) * A) /\ IsHermiteNormalForm ((V : Mat) * B) /\
    (U : Mat) * A != (V : Mat) * B /\
    (forall H1 H2, RowEquiv A H1 -> RowEquiv B H2 -> H1 != H2) /\
    (forall H1 H2, RowEquiv A H1 -> RowEquiv B H2 ->
        IsEmpty (coker H1 ~+ coker H2)) /\
    (forall P Q : GL (Fin 2) Int, (P : Mat) * B * Q != A) /\
    exists Y : Mat, (Y * U * V).det = (Y * V * U).det /\
      (forall H1 H2, RowEquiv (Y*U*V) H1 -> RowEquiv (Y*V*U) H2 -> H1 != H2) /\
      (exists H1 H2, RowEquiv (Y*U*V) H1 /\ IsHermiteNormalForm H1 /\
                     RowEquiv (Y*V*U) H2 /\ IsHermiteNormalForm H2)
\end{verbatim}
(ASCII transcription; the source uses the Unicode symbols, \texttt{Mat} is
\texttt{Matrix (Fin 2) (Fin 2) Int}, \texttt{\~{}+} is Mathlib's additive equivalence
\texttt{AddEquiv}, and $Y\,U\,V$ is the product of matrices with $U,V$ coerced from
$\GL$.) Correspondence: \texttt{det} is Mathlib's \texttt{Matrix.det}; unimodular
transformations are Mathlib's \texttt{GL (Fin 2) Int}, with noncommutativity stated in
the group; \texttt{RowEquiv A H} is $\exists W\in\GL_2(\Z),\ WA=H$;
\texttt{coker X} is the quotient module $\Z^2/L(X)$ with \texttt{rowModule X} $=L(X)$
and \texttt{rowModule\_eq\_span} identifying it with the span of the rows. The witnesses
are $A=M$, $B=N$, $Y=X$ and the pair $U,V$ above.

Checks: \texttt{lake build} succeeds, \texttt{lake env lean -DwarningAsError=true} on
the file succeeds, there is no \texttt{sorry} or \texttt{native\_decide}, and
\texttt{\#print axioms Conjecture5600.conjecture\_5600} reports only
\texttt{propext}, \texttt{Classical.choice}, \texttt{Quot.sound}.

\section*{Credit}
The pair $\operatorname{diag}(1,4)$, $\operatorname{diag}(2,2)$ and the parity
obstruction in part (3) are adapted from the accepted solution of the near-duplicate
conjecture 00000006330 by gaochengzhecpu
(\url{solutions/00000006330/gaochengzhecpu_submission_20261003215646/}), used under the
repository's GPL-3.0 license. The Mathlib formalization, the cokernel invariant, the
unimodular pair and reading (B) are new here.

\end{document}
